\scanpage{126}
% line head | text | 75
\mspage{75}
% line 1–3 | text | Prop 4
\uline{Prop 4} : Soit $t\in\inserted{\operatorname{Ob}}T$. Pour que $t\in\inserted{\operatorname{Ob}}\bar U$, il f. et s. que
$q(t)\in\inserted{\operatorname{Ob}}\overline{U'}$. En d'autres termes, le ``bord'' $\dot T=T\cap\bar U$ de $T$
dans $X$ est ``le même'' que son bord dans $X'$.

% margin note written sideways
\margin{\uline{Proposition 5} : Soit
\diagram{\begin{tikzcd}[ampersand replacement=\&,row sep=tiny,column sep=tiny]
Y\arrow[r,"i"]\arrow[d,"p"'] \& X\arrow[d,"q"]\\
Y'\arrow[r,"i'"'] \& X'
\end{tikzcd}}
un \deleted{\unsure{ca}}carré cocartésien dans $\mathrm{Cat}$. Si $i$ est une immersion ouverte (resp. fermée) alors $i'$ aussi, \inserted{et le carré est aussi cartésien.} \unsure{idem|id.} pour $p$ et $q$. Si $i$ et $p$ sont des imm. ouvertes \inserted{(resp. \unsure{fermées})} \deleted{\unsure{alors}} la situation est isom. à celle d'un objet $X'$ de $\mathrm{Cat}$, de deux ouverts (resp. fermés) $X,Y'$ de $X$ ay. \unsure{pr|pour} \unsure{réunion|somme} $X'$, pour intersection $Y$. Inversement, dans une telle situation, \deleted{$X$ s'id} le carré correspondant est \unsure{cocartésien|cartésien}, i.e. $X$ s'identifie à la somme amalgamée
$X'\simeq X\amalg_Y Y'$.
\unsure{Nous retrouvons|Pour ceci cf} \unsure{en effet|au §} la prop 15, page 158, et son cor. \unsure{Ils disent}}
% line 4 | text | Dém. Dire que
Dém. Dire que $t\in\operatorname{Ob}\bar U$ signifie aussi
% line 5 | display | U over t nonempty
\[
U_{/t}\ne\varnothing
\]
% line 6–8 | text | et implique, par continuité
et implique, par continuité, que $q(t)\in\deleted{\unsure{\operatorname{Ob}}}\inserted{\overline{q(U)}\subset}\overline{U'}$. Inversement,
si $q(t)\in\overline{U'}$, i.e. s'il existe $u'\in U'$
tel que
% line 9 | display | Hom(u prime,q(t)) nonempty
\[
\deleted{\operatorname{Hom}(u'}\quad\operatorname{Hom}(u',q(t))
\bigl(=\pi_0({}_{u'}\backslash^{U}_{/t})\bigr)\ne\varnothing,
\]
% line 10–11 | text | en résulte que
en résulte que $\mathord{}_{u'}\backslash^{U}_{/t}\ne\varnothing$, a fortiori
$U_{/t}$ (qui \unsure{reçoit} $\mathord{}_{u'}\backslash^{U}_{/t}$) $\ne\varnothing$, \unsure{qed}
\margin{\uline{NB} Nous utilisons ``\uline{immersion ouverte}'', ``\uline{immersion fermée}'' au \uline{sens strict} : c'est un \uline{monomorphisme} \unsure{pleinement|entièrement} fidèle dont l'image \deleted{\unsure{essentielle}} est une sous-catégorie \uline{pleine} \unsure{strictement} ouverte ou fermée.}

% line 12 | text | 2. Sommes amalgamées et W-équivalences
2\unsure{.}\quad \uline{Sommes amalgamées et $W$-équivalences}.

% line 13 | text | 2.1. Axiome du carré cocartésien
2.1. \uline{Axiome du carré cocartésien et axiome du cube} (ou ``\uline{lemme des cinq}'')

% line 14–15 | text | Nous supposons à présent donné
Nous supposons à présent donné un localiseur
fondamental
% line 16 | display | W contained in Fl(Cat)
\[
W\subset\operatorname{Fl}(\mathrm{Cat}),
\]
% line 17–24 | text | qui satisfait à W(1,2bis,3)
qui satisfait à $W(1,2\mathrm{bis},3)$. Nous ne supposons
pas pour le moment \unsure{pour|que} les \deleted{\unsure{autres}} autres axiomes
\unsure{suivants|qu'il ait} $W(4.)$ \deleted{$\unsure{W(5)}$} (axiome de fibration), $W(5)$
(axiome de \deleted{limites inductives}\inserted{\unsure{connexité}}), $W(6)$ (axiome des
limites inductives). Mais nous allons introduire
à présent deux axiomes \inserted{de stabilité} supplémentaires, de
nature différente, \deleted{\unsure{car ils}} concernant les sommes
amalgamées.

% line 25–28 | text | W(7bis), Axiome du carré cocartésien
\uuline{$W(7\mathrm{bis})$} \inserted{(Axiome du carré cocartésien)} Soient $X$ un objet de $\mathrm{Cat}$, $U,U'$ deux
ouverts de $X$ tel que $U\cup U'=X$, considérons le
diagramme cocartésien correspondant (cf prop. 5 en
\deleted{\unsure{marge}} de la page) :
% archive sheet number: 126
